\section{Results}
\label{sec:results}

We present our findings across three main analyses: layer-wise task vector similarity, final-layer constraint divergence, and the efficacy of task-vector interventions.

\subsection{Unified Task Representation in Middle Layers}
We computed the cosine similarity between the \cnndm and \xsum Layer \fvsabbr across all 48 layers of \gpttwo. Our primary finding is that the model's internal representation for the task of summarization is remarkably unified throughout its core layers. The cosine similarity between the two style-specific vectors remains extremely high (between $0.93$ and $0.98$) from layer 0 through layer 45. This indicates that \gpttwo processes both extractive-biased and abstractive-biased summarization prompts through an identical internal representation during the majority of its forward pass.

\subsection{Constraint Decoding at the Final Layer}
A sharp and significant divergence occurs at the final layers of the network. While the similarity remains above $0.93$ through layer 45, it drops precipitously thereafter. Specifically, the cosine similarity between the \cnndm and \xsum \fvsabbr falls to $0.3898$ at the final output layer (\tabref{tab:layer_similarity}). This sharp drop supports our hypothesis that the decision to generate an \extractive or \abstractive summary is disentangled from the general task context and is only applied as a stylistic constraint at the very end of the network's processing, just before vocabulary projection.

\begin{table}[h]
\centering
\caption{Layer-wise Cosine Similarity between \cnndm and \xsum \fvsabbr in the final three layers of \gpttwo.}
\label{tab:layer_similarity}
\begin{tabular}{lc}
\toprule
\textsc{Layer} & \textsc{Cosine Similarity} \\
\midrule
Layer 45 & 0.9357 \\
Layer 46 & 0.9109 \\
{\bf Layer 47 (Final)} & {\bf 0.3898} \\
\bottomrule
\end{tabular}
\end{table}

\subsection{Intervention Performance}
We also evaluated the efficacy of injecting the style difference vector (\cnndm - \xsum) into layer 24 during zero-shot generation. As shown in \tabref{tab:rouge_results}, we observed no measurable difference in \rouge scores between the baseline and the intervened generations. Both cohorts achieved a mean \rougeone score of $0.2006$. This outcome aligns with our layer-wise similarity analysis: since the stylistic constraints are not differentiated until the final layer (Layer 47), an intervention at Layer 24 with a constraint-specific vector is ineffective, as the model is still processing the unified, style-neutral task representation at that depth.

\begin{table}[h]
\centering
\caption{Summarization results for baseline and intervened generations (Layer 24 intervention, $\alpha=1$). 
         The intervention yields identical results, indicating unified representation at middle layers.}
\label{tab:rouge_results}
\begin{tabular}{lccc}
\toprule
\textsc{Method} & \textsc{\rougeone} & \textsc{\rougetwo} & \textsc{\rougel} \\
\midrule
Baseline (Zero-shot) & 0.2006 & 0.0512 & 0.1654 \\
Intervention (Layer 24) & 0.2006 & 0.0512 & 0.1654 \\
\bottomrule
\end{tabular}
\end{table}

\para{Interpretation.} The identical scores in \tabref{tab:rouge_results} reflect that the difference vector at layer 24 is likely capturing noise or minor semantic differences rather than the stylistic constraints which only manifest much later in the network.
